\chapter{Contrôles de risque}\label{ch:hf:risk-controls}

Le 1\up{er} août 2012, le système de routage d'un teneur de marché a envoyé des millions d'ordres sur le marché en environ quarante-cinq minutes. Il a obtenu plus de 4
millions d'exécutions sur 154 actions, pour plus de 397 millions d'actions, et perdu plus de 460~millions de dollars. L'ordonnance du régulateur, rendue ensuite, se lit comme une
liste des contrôles qui lui manquaient. Dans la reconstruction de ce chapitre, une \omterm{def:hf:risk-controls:loss}{limite de perte} de 2~millions de dollars, vérifiée chaque seconde, aurait arrêté
la même course folle au bout de douze secondes et de 2.1~millions de dollars, sans se déclencher une seule fois sur 2\,500 jours ordinaires~; un limiteur de débit d'ordres, le contrôle le plus
souvent cité en premier, ne l'aurait pas arrêtée du tout.

\section{Limites pré-négociation}

Les ordres d'un teneur de marché passent par une porte de contrôle pré-négociation avant d'atteindre une plateforme~: un ensemble de vérifications qui refusent tout ordre enfreignant une limite. Le
livre~13, chapitre~22, construit la porte qui exécute ces vérifications en nanosecondes~; ce chapitre fixe la politique qu'elle applique.

\begin{definition}[Contrôle de risque pré-négociation]\label{def:hf:risk-controls:pretrade}
Un \emph{contrôle de risque pré-négociation}\index{contrôle de risque pré-négociation} est un test qu'un ordre doit réussir avant d'être envoyé à une plateforme, au regard de limites portant sur sa taille,
sa valeur et son prix, ainsi que sur les positions, les ordres ouverts et les expositions qu'il créerait, l'ordre étant refusé si une limite serait enfreinte.
\end{definition}

\begin{definition}[Tunnel de prix]\label{def:hf:risk-controls:collar}
Un \emph{tunnel de prix}\index{tunnel de prix} (price collar) est une limite pré-négociation sur l'écart maximal entre le prix d'un ordre et un prix de référence (la dernière transaction, le mid, une
fair value), exprimé en points de base ou en ticks~; au-delà duquel un achat, ou en deçà duquel une vente, est refusé comme erreur probable.
\end{definition}

Les limites forment une hiérarchie (voir la \cref{fig:hf:risk-controls:hierarchy})~: la société, ses desks, leurs stratégies, puis les instruments. Chaque niveau plafonne ce que
peuvent cumuler les niveaux inférieurs. Le schéma du chapitre (\texttt{firm.riskctl}) l'écrit sous la forme d'un instantané que la porte peut charger~: exposition brute et \omterm{def:hf:risk-controls:loss}{limites de perte}
pour la société et pour chaque desk, un débit d'ordres, des ordres ouverts, une \omterm{def:hf:risk-controls:loss}{limite de perte} et une limite de position pour chaque stratégie, ainsi qu'un collar, une quantité maximale
un notionnel maximal et des limites à l'achat et à la vente pour chaque instrument. Un instantané est validé avant emploi (chaque desk dans les limites de la société, chaque
stratégie rattachée à un desk connu) puis exporté avec la section de la porte, que la porte du livre~13 lit telle quelle.

\begin{omfigure}
\begin{tikzpicture}[font=\footnotesize,>=stealth,box/.style={draw,rounded corners,fill=omThm!8,minimum width=2.4cm,minimum height=0.7cm,align=center}]
\node[box] (f) at (0,2.4) {société\\exposition brute, perte, messages};
\node[box] (d1) at (-2.2,1.2) {desk\\exposition brute, perte};
\node[box] (d2) at (2.2,1.2) {desk\\exposition brute, perte};
\node[box] (s1) at (-3.9,0) {stratégie\\débit, ouverts, perte};
\node[box] (s2) at (-0.9,0) {stratégie};
\node[box] (s3) at (2.2,0) {stratégie};
\node[box,fill=gray!10] (i) at (0,-1.2) {instruments~: collar, taille, notionnel, longs, shorts};
\draw[->] (f) -- (d1); \draw[->] (f) -- (d2); \draw[->] (d1) -- (s1); \draw[->] (d1) -- (s2); \draw[->] (d2) -- (s3);
\node[box,fill=omDef!10] (g) at (6.8,1.2) {porte pré-négociation\\(refus)};
\node[box,fill=omDef!10] (m) at (6.8,0) {moniteur post-négociation\\(perte, position, débit)};
\node[box,fill=omDef!20] (k) at (6.8,-1.2) {kill switch\\(annulation, aplatissement, blocage)};
\draw[->] (m) -- (k);
\draw[->,dashed] (2.2,2.4) -- node[above] {instantané des limites} (6.8,2.4) -- (g);
\end{tikzpicture}

\omcaption{La hiérarchie des limites et les trois endroits où elle agit~: la porte pré-négociation refuse un ordre qui enfreint une limite~; le moniteur post-négociation surveille
les pertes, les positions et les débits de messages et déclenche le \omterm{def:hf:risk-controls:kill}{kill switch}, qui annule les ordres ouverts, aplatit les positions et bloque les nouveaux ordres au niveau
qui a été franchi. Source~: \texttt{firm.riskctl}~; la porte est le \texttt{firm.riskgate} du livre~13.}\label{fig:hf:risk-controls:hierarchy}
\end{omfigure}

\section{Limites de position, de perte et de messages}

\begin{definition}[Limite de perte]\label{def:hf:risk-controls:loss}
Une \emph{limite de perte}\index{limite de perte} est un seuil portant sur la perte réalisée augmentée de la perte en valeur de marché d'une stratégie, d'un desk ou de la société sur une période (une journée,
en général), au-delà duquel la perte déclenche une action d'arrêt à ce niveau.
\end{definition}

Les limites de position et de perte sont vérifiées après les transactions comme avant elles~: la porte peut refuser un ordre qui dépasserait une position, mais seul un
moniteur qui valorise les positions au marché voit une perte se constituer (voir le \cref{lst:hf:risk-controls:monitor}). Une \omterm{def:hf:risk-controls:loss}{limite de perte} est le contrôle le plus direct de tous,
puisqu'elle borne ce qu'elle protège, mais elle n'agit qu'une fois la perte survenue et seulement aussi souvent que les positions sont valorisées. Les limites de messages
(les limiteurs et les ratios ordres sur transactions du chapitre~10) protègent la plateforme et la réputation de la société auprès d'elle~; la limite que la société fixe au nombre de messages par seconde est un
moniteur de l'ensemble de la société, et non d'une stratégie.

\section{Tunnels de prix et contrôles de saisie}

Les tunnels et les contrôles de taille repèrent l'erreur dans un ordre unique~: un prix loin du marché, une quantité ou un notionnel très au-dessus de la normale. Le 2 mai 2022, une
banque, un trader qui voulait vendre un panier de 58~millions de dollars en a créé un de 444~milliards de dollars~; les contrôles de la banque en ont bloqué 255~milliards de dollars, mais en ont laissé passer
189~milliards de dollars jusqu'à un algorithme de trading, qui a vendu environ 1.4~milliard de dollars sur les marchés européens avant d'être annulé. Le régulateur britannique de la conduite des affaires, qui
a infligé à la banque une amende de 27.8~millions de livres en 2024, a constaté qu'aucun blocage ferme n'aurait rejeté le panier dans sa totalité, que le trader pouvait fermer une alerte surgissante
sans la lire, et que la surveillance en temps réel était trop lente.

\begin{dated}{2026-09}{Les règles sur l'accès au marché et le trading algorithmique}\label{dat:hf:risk-controls:rules}
La SEC a adopté la Rule 15c3-5 en novembre 2010~: les courtiers ou dealers ayant un accès au marché doivent maintenir des contrôles de gestion des risques raisonnablement conçus
pour empêcher les ordres erronés et les ordres qui dépassent les seuils de crédit et de capital fixés à l'avance, et leur directeur général doit les certifier. Son ordonnance du
16 octobre 2013 a infligé au teneur de marché de la course folle de 2012 une amende de 12~millions de dollars pour violation de la règle. Dans l'Union européenne, le Commission Delegated
Regulation (EU) 2017/589 exige qu'une entreprise d'investissement puisse annuler immédiatement, à titre de mesure d'urgence, tout ou partie de ses ordres non exécutés
sur tout ou partie des plateformes, et identifier l'algorithme et le trader responsables de chaque ordre. La Financial Conduct Authority britannique a infligé à une banque
une amende de 27\,766\,200~livres le 22 mai 2024 pour l'erreur de panier de 2022, et la Prudential Regulation Authority une amende supplémentaire de 33\,880\,000~livres.
\end{dated}

\begin{definition}[Market Access Rule]\label{def:hf:risk-controls:access}
La \emph{Market Access Rule}\index{Market Access Rule} est la Rule 15c3-5 de la SEC, qui exige qu'un broker-dealer ayant accès à une bourse ou à un système de négociation
alternatif (ATS), pour son propre compte ou pour celui de ses clients, dispose de contrôles pré-négociation documentés et régulièrement revus contre les ordres erronés et les ordres
au-delà des seuils de crédit et de capital fixés à l'avance, sous son contrôle direct et exclusif.
\end{definition}

\section{Kill switches~: qui les actionne et quand}

\begin{definition}[Kill switch]\label{def:hf:risk-controls:kill}
Un \emph{kill switch}\index{kill switch} est un contrôle qui, lorsqu'il est déclenché par une personne ou par un moniteur, arrête la négociation d'une stratégie, d'un desk ou de la société entière
d'un seul coup~: il annule ses ordres ouverts, refuse les nouveaux et, si cela est configuré ainsi, aplatit ses positions, et il enregistre qui l'a déclenché et pourquoi.
\end{definition}

La course folle du chapitre (voir le \cref{lst:hf:risk-controls:runaway}) reprend les chiffres de l'ordonnance de la SEC~: environ 1\,480 ordres par seconde de 100 actions pendant 45
minutes, une perte de 1.16~\$ par action négociée, et 2.46~millions de dollars par seconde de position brute (les 3.5~milliards de dollars à l'achat et les 3.15~milliards de dollars à la vente où elle
s'est achevée). L'arrêt aplatit la position à 0.2~\% de sa valeur. Chaque contrôle, seul, l'arrête à un moment différent
(voir la \cref{fig:hf:risk-controls:controls}), et chacun se déclenche sur une part des jours ordinaires d'un livre légitime (débits de pointe d'environ 300 messages
par seconde, positions brutes d'environ 150~millions de dollars, drawdowns intrajournaliers d'environ 150\,000~\$, ordres à quelques dizaines de points de base de leurs références).

\begin{center}
\small
\setlength{\tabcolsep}{4pt}%
\begin{tabular}{@{}>{\raggedright\arraybackslash}p{6.0cm}>{\centering\arraybackslash}p{2.6cm}>{\centering\arraybackslash}p{2.6cm}>{\centering\arraybackslash}p{2.6cm}@{}}
\toprule
contrôle & arrêt après & perte & se déclenche les jours ordinaires\\
\midrule
aucun (les 45 minutes de 2012) & 2\,700 s & 477~millions de dollars & --\\
une personne à 5 minutes & 300 s & 53~millions de dollars & --\\
limiteur de débit, 500 par seconde & jamais & 161~millions de dollars & 10.2~\%\\
\omterm{def:hf:risk-controls:collar}{tunnel de prix}, 5~\% & 600 s & 106~millions de dollars & 0~\%\\
seuil de capital, 1~milliard de dollars brut & 407 s & 72~millions de dollars & 0~\%\\
limite de position, 250~millions de dollars brut & 102 s & 18~millions de dollars & 7.6~\%\\
\omterm{def:hf:risk-controls:loss}{limite de perte}, 2~millions de dollars, chaque seconde & 12 s & 2.1~millions de dollars & 0~\%\\
tous ensemble & 35 s & 2.1~millions de dollars & 10.2~\%\\
\bottomrule
\end{tabular}
\end{center}

\begin{omfigure}
\begin{tikzpicture}
\begin{axis}[width=11cm,height=5.4cm,ybar,bar width=12pt,ymode=log,log basis y=10,ymin=1,ymax=900,xmin=-0.6,xmax=7.6,
  xtick={0,1,2,3,4,5,6,7},xticklabels={aucun,{personne, 5 min},limiteur,tunnel,capital,position,{limite de perte},tous},
  xticklabel style={font=\scriptsize,rotate=25,anchor=north east},ylabel={\small perte de l'emballement (millions de dollars)},
  tick label style={font=\footnotesize},grid=major,grid style={gray!25},table/col sep=comma,nodes near coords,
  nodes near coords style={font=\scriptsize},point meta=rawy]
\addplot[fill=omThm!60,draw=omThm] table[x=i,y=loss]{figdata/hft/27-risk-controls/controls.csv};
\end{axis}
\end{tikzpicture}

\omcaption{La perte d'un emballement, calquée sur l'incident de 2012, selon chaque contrôle seul puis selon l'ensemble de ces contrôles (échelle logarithmique, en millions de dollars)~:
aucun contrôle (arrêté à la main au bout de 45 minutes), une personne à 5 minutes, un limiteur de débit de 500 ordres par seconde, un \omterm{def:hf:risk-controls:collar}{tunnel de prix} de 5~\%, un seuil de capital
de 1~milliard de dollars, une limite de position de 250~millions de dollars et une \omterm{def:hf:risk-controls:loss}{limite de perte} de 2~millions de dollars contrôlée chaque seconde. Données~: \texttt{hf\_risk.table}.}\label{fig:hf:risk-controls:controls}
\end{omfigure}

Le limiteur de débit, souvent le premier contrôle cité, n'arrête rien~: il ralentit l'emballement jusqu'au tiers, et la perte suit, soit 161~millions de dollars à la fin.
Le \omterm{def:hf:risk-controls:collar}{tunnel de prix} et le seuil de capital, les contrôles que la \omterm{def:hf:risk-controls:access}{Market Access Rule} impose, l'arrêtent au bout de dix et sept minutes, pour 106 et 72
millions de dollars. La limite de position l'arrête en moins de deux minutes~; la \omterm{def:hf:risk-controls:loss}{limite de perte}, en douze secondes. Ensemble, ils l'arrêtent à 35 secondes~: le limiteur de débit,
en ralentissant l'emballement, retarde la \omterm{def:hf:risk-controls:loss}{limite de perte}, mais la perte reste la même.

Chaque seuil est un compromis (la \cref{fig:hf:risk-controls:tradeoff})~: une \omterm{def:hf:risk-controls:loss}{limite de perte} plus serrée arrête l'emballement plus tôt, mais se déclenche les jours ordinaires, dans 2~\%
des cas à 1~million de dollars, dans 19~\% des cas à 500\,000~dollars~; une limite de position assez serrée pour l'arrêter en une minute se déclenche sur la moitié des jours. La \omterm{def:hf:risk-controls:loss}{limite de perte} domine~:
pour le même nombre de fausses alertes, elle arrête l'emballement bien plus tôt, parce qu'elle mesure ce qui compte. La limite de position attrape ce que la \omterm{def:hf:risk-controls:loss}{limite de perte}
ne peut pas voir~: un emballement qui n'a pas encore perdu d'argent.

\begin{omfigure}
\begin{tikzpicture}
\begin{axis}[width=11cm,height=5.4cm,xlabel={\small jours ordinaires où le contrôle se déclenche (\%)},ylabel={\small perte de l'emballement (millions de dollars)},
  xmin=-2,xmax=92,ymode=log,log basis y=10,ymin=0.4,ymax=60,tick label style={font=\footnotesize},grid=major,grid style={gray!25},
  table/col sep=comma,legend style={font=\scriptsize,at={(0.97,0.97)},anchor=north east}]
\addplot[omThm,thick,mark=*,mark size=1.4pt] table[x=false,y=loss]{figdata/hft/27-risk-controls/sweep_loss.csv};
\addplot[omDef,thick,dashed,mark=square*,mark size=1.3pt] table[x=false,y=loss]{figdata/hft/27-risk-controls/sweep_position.csv};
\legend{{limite de perte, 0.5--4 millions de dollars},{limite de position, 100--600 millions de dollars}}
\end{axis}
\end{tikzpicture}

\omcaption{Le compromis d'un seuil~: la perte de l'emballement sous une \omterm{def:hf:risk-controls:loss}{limite de perte} (de 0.5 à 4~millions de dollars) et sous une limite de position brute (de 100 à
600~millions de dollars), en regard de la part des 2\,500 jours ordinaires où chacune se serait déclenchée. Données~: \texttt{hf\_risk.sweep}.}\label{fig:hf:risk-controls:tradeoff}
\end{omfigure}

Qui actionne le \omterm{def:hf:risk-controls:kill}{kill switch} compte autant que le moment. Un moniteur réagit en une seconde~; une personne met des minutes à voir, comprendre et agir, et en 2012 elle a mis
quarante-cinq minutes. Le \omterm{def:hf:risk-controls:kill}{kill switch} doit être accessible au moniteur comme à une personne, doit annuler avant de liquider, doit agir sur toutes les plateformes à la fois (comme
l'exige la règle européenne), et doit laisser une trace d'audit de qui l'a actionné et pourquoi (\texttt{firm.riskctl.KillSwitch}).

\section{Exigences réglementaires pour les firmes algorithmiques}

Les règles de l'encadré daté convergent vers les mêmes contrôles~: limites avant négociation sur le prix, la taille et la valeur~; limites de crédit et de capital~; une fonction
de coupure~; savoir quel algorithme et quel trader ont envoyé chaque ordre~; et une revue des contrôles par la direction. Les injonctions des régulateurs se lisent comme des descriptions
de ce qui a mal tourné lorsqu'un de ces contrôles manquait~: en 2012, aucun contrôle du capital global engagé par les ordres, et aucune alerte automatisée n'a conduit à agir
à temps~; en 2022, un contrôle de taille pouvait être contourné et la surveillance était trop lente pour agir.

\section{Tutoriel~: quarante-cinq minutes}\label{tut:hf:risk-controls:tut}

\begin{tutorial}
\textbf{Objectif.} Écrire l'instantané des limites, l'exporter pour la passerelle, et mesurer ce que chaque contrôle aurait sauvé face à un emballement de type 2012 et
ce qu'il coûte les jours ordinaires. \textbf{État final~:} le tableau et les trois figures.
\begin{tutsteps}
\item \textbf{L'instantané}~: \texttt{firm.riskctl.validate} et \texttt{export}~; \texttt{riskctl\_fixture.py} écrit
\texttt{data/limits.json}, que charge \texttt{firm.riskgate.from\_riskctl} du livre~13.
\item \textbf{Le moniteur}~: les marques, les \omterm{def:hf:risk-controls:loss}{limites de perte} et les actions de coupure à chaque niveau.
\omcode{code/firm/riskctl/firm_riskctl.py}{126}{143}{Met les positions au prix du marché, somme les stratégies par desk et pour la firme, et coupe au niveau dont la limite de perte est franchie.}\label{lst:hf:risk-controls:monitor}
\item \textbf{L'emballement} sous chaque contrôle.
\omcode{code/firm/riskctl/firm_riskctl.py}{169}{188}{Le temps d'arrêt de chaque contrôle~; le premier qui arrête l'emballement, lequel est alors liquidé.}\label{lst:hf:risk-controls:runaway}
\item \textbf{Les jours ordinaires} et les balayages de seuils (\texttt{hf\_risk.table}, \texttt{hf\_risk.sweep}).
\end{tutsteps}
\textbf{Pour aller plus loin.} Rejouer l'emballement à travers la passerelle du livre~13 sur \texttt{firm.exchsim} (sa \texttt{firm\_riskgate\_runaway}), ajouter un
second emballement qui perd lentement, et voir quel contrôle le détecte~; rendre la \omterm{def:hf:risk-controls:loss}{limite de perte} fondée sur le drawdown intrajournalier.
\end{tutorial}

\section{Construction~: les contrôles de risque}\label{bld:hf:risk-controls:riskctl}

\begin{build}
\bfield{Rôle} Assurer la politique des limites (hiérarchie, schéma, validation), les moniteurs post-négociation et le \omterm{def:hf:risk-controls:kill}{kill switch}~; fournir à la passerelle pré-négociation
son instantané.
\bfield{Interface} \texttt{validate(limits)}, \texttt{to\_riskgate(limits)}, \texttt{export(limits)}, \texttt{KillSwitch},
\texttt{Monitor(limits, switch)} avec \texttt{on\_fill}, \texttt{mark}, \texttt{on\_message}~; \texttt{runaway(controls)}, \texttt{normal\_days(controls)}. Le schéma,
version 1, figure dans la docstring du module et constitue le contrat avec \texttt{firm.riskgate} du livre~13.
\bfield{Règles} Dollars, actions, nanosecondes~; un desk dans les limites de la firme~; les coupures enregistrées avec leur motif~; annuler avant de liquider.
\bfield{Tests d'acceptation} \texttt{code/firm/riskctl/tests/}~: la validation détecte un desk au-dessus de la firme et un desk inconnu~; la passerelle de la section à la main~;
l'instantané exporté est identique au fixture, et la passerelle du livre~13 le charge et refuse un ordre hors du collar~; le moniteur coupe une stratégie au-delà de sa \omterm{def:hf:risk-controls:loss}{limite de perte}
et la firme au-delà de son débit de messages, et l'audit enregistre la coupure et la reprise~; l'emballement et les jours ordinaires se comportent comme dans le chapitre.
\bfield{Pour aller plus loin} Limites de drawdown intrajournalier~; coupure par plateforme~; limites modifiées au cours de la journée avec des instantanés versionnés.
\end{build}

\begin{omsources}
\item US Securities and Exchange Commission, \emph{In the Matter of Knight Capital Americas LLC}, Release No.\ 34-70694, 16 octobre 2013.
\item Commission Delegated Regulation (EU) 2017/589, article 12 (fonction de coupure).
\item Financial Conduct Authority, FCA fines CGML £27\,766\,200 for failures in its trading systems and controls, 22 mai 2024.
\end{omsources}

\section{Exercices}

\begin{exercise}[$\star$]\label{exo:hf:risk-controls:1}
L'emballement perd 1.16~\$ par action sur 148\,100 actions par seconde. Au bout de combien de temps une \omterm{def:hf:risk-controls:loss}{limite de perte} de 2~millions de dollars, contrôlée chaque seconde, se déclenche-t-elle~?
\end{exercise}

\begin{exercise}[$\star$]\label{exo:hf:risk-controls:2}
La position brute croît de 2.46~millions de dollars par seconde. Quand une limite de position de 250~millions de dollars l'arrête-t-elle, et un seuil de capital de 1~milliard de dollars~?
\end{exercise}

\begin{exercise}[$\star$]\label{exo:hf:risk-controls:3}
Un ordre d'achat est fixé à 42.10~\$ pour une référence de 40.00~\$. Un \omterm{def:hf:risk-controls:collar}{tunnel de prix} de 5~\% le refuse-t-il~?
\end{exercise}

\begin{exercise}[$\star\star$]\label{exo:hf:risk-controls:4}
Pourquoi le limiteur de débit n'arrête-t-il pas l'emballement, et à quoi sert-il~?
\end{exercise}

\begin{exercise}[$\star\star$]\label{exo:hf:risk-controls:5}
Pourquoi ajouter tous les contrôles arrête-t-il l'emballement plus tard (35 secondes) que la seule \omterm{def:hf:risk-controls:loss}{limite de perte} (12)~?
\end{exercise}

\begin{exercise}[$\star\star$]\label{exo:hf:risk-controls:6}
Quelle doit être la première action d'un \omterm{def:hf:risk-controls:kill}{kill switch}~: annuler les ordres ou liquider les positions~? Pourquoi~?
\end{exercise}

\begin{exercise}[$\star\star\star$]\label{exo:hf:risk-controls:7}
\emph{Programmation.} Avec \texttt{hf\_risk.sweep}, quelle \omterm{def:hf:risk-controls:loss}{limite de perte} se déclenche sur au plus 1~\% des jours ordinaires, et que laisse perdre l'emballement~?
\end{exercise}

\begin{exercise}[$\star\star\star$]\label{exo:hf:risk-controls:8}
\emph{Trouver la faille.} «~Notre algorithme n'a jamais perdu plus de 500\,000~\$ en une journée, donc une \omterm{def:hf:risk-controls:loss}{limite de perte} de 500\,000~\$ ne nous coûte rien.~»
\end{exercise}

\section{Problème~: quarante-cinq minutes}

\begin{problem}[{Problème du week-end --- quarante-cinq minutes}]\label{pb:hf:risk-controls:1}
Une entreprise de teneur de marché fixe ses contrôles de risque après avoir lu l'ordonnance de 2012.

\medskip\noindent\textbf{Partie I --- L'incident et les règles.}
\begin{enumerate}
\item Résumer ce qui s'est passé le 1\up{er} août 2012, selon l'ordonnance de la SEC.
\item Définir la \omterm{def:hf:risk-controls:access}{Market Access Rule} et ce qu'elle exige.
\item Que exige la fonctionnalité de kill de l'UE~?
\item Qu'est-ce qui a mal tourné dans l'erreur de panier de 2022, selon le compte rendu de la FCA~?
\end{enumerate}

\medskip\noindent\textbf{Partie II --- Les contrôles.}
\begin{enumerate}[resume]
\item Définir un \omterm{def:hf:risk-controls:pretrade}{contrôle de risque pré-négociation} et un \omterm{def:hf:risk-controls:collar}{tunnel de prix}.
\item Définir une \omterm{def:hf:risk-controls:loss}{limite de perte} et un \omterm{def:hf:risk-controls:kill}{kill switch}.
\item Décrire la hiérarchie des limites et la validation de l'instantané.
\item Comment l'instantané atteint-il la porte de la nanoseconde~?
\end{enumerate}

\medskip\noindent\textbf{Partie III --- Mesures.}
\begin{enumerate}[resume]
\item Donner le temps d'arrêt et la perte sous chaque contrôle.
\item Quels contrôles se déclenchent les jours ordinaires, et à quelle fréquence~?
\item Décrire les compromis entre \omterm{def:hf:risk-controls:loss}{limite de perte} et limite de position.
\item Pourquoi la \omterm{def:hf:risk-controls:loss}{limite de perte} domine-t-elle, et qu'apporte la limite de position~?
\end{enumerate}

\medskip\noindent\textbf{Partie IV --- Le verdict.}
\begin{enumerate}[resume]
\item Énoncer le \emph{résultat nommé}~: la perte que l'emballement aurait subie sous chaque contrôle, et le premier contrôle qui l'arrête.
\item Quel ensemble de contrôles feriez-vous fonctionner, avec quels seuils~?
\item Qui doit pouvoir actionner le \omterm{def:hf:risk-controls:kill}{kill switch}, et en combien de temps~?
\item Que doit consigner la piste d'audit~?
\item De quoi a besoin un emballement qui perd lentement~?
\item Comment tester les contrôles en l'absence d'emballement~?
\item Que doit certifier le directeur général d'une entreprise~?
\item En une phrase~: que rapporte un contrôle de risque~?
\end{enumerate}
\end{problem}

\section{Questions d'entretien}

\begin{interviewq}[$\star$ \iqroles{risque}]\label{iq:hf:risk-controls:1}
Énumérez les contrôles pré-négociation que vous exigeriez sur chaque ordre, dans l'ordre où vous les exécuteriez.
\end{interviewq}

\begin{interviewq}[$\star\star$ \iqroles{développeur}]\label{iq:hf:risk-controls:2}
Votre passerelle refuse les ordres lorsque son instantané des limites est périmé. Pourquoi échouer en position fermée, et que pourrait-il mal se passer~?
\end{interviewq}

\begin{interviewq}[$\star\star$ \iqroles{trader}]\label{iq:hf:risk-controls:3}
Votre stratégie a atteint sa \omterm{def:hf:risk-controls:loss}{limite de perte} à 10~h~05 un jour volatil et a été liquidée. Le marché s'est redressé à 10~h~30. La limite était-elle fausse~?
\end{interviewq}

\begin{interviewq}[$\star\star$ \iqroles{risque}]\label{iq:hf:risk-controls:4}
Comment fixeriez-vous la limite de position d'une stratégie à partir de son historique~?
\end{interviewq}

\begin{interviewq}[$\star\star$ \iqroles{développeur}]\label{iq:hf:risk-controls:5}
Concevez le \omterm{def:hf:risk-controls:kill}{kill switch} d'une entreprise connectée à 20 plateformes par l'intermédiaire de 3 passerelles. Que doit-il se produire dans la première milliseconde~?
\end{interviewq}

\begin{interviewq}[$\star\star\star$ \iqroles{chercheur}]\label{iq:hf:risk-controls:6}
Les replis quotidiens ordinaires suivant une loi normale de moyenne $\mu$ et d'écart type $\sigma$, et un emballement perdant à la vitesse $\ell$ par seconde, déterminez la \omterm{def:hf:risk-controls:loss}{limite de perte}
qui minimise le coût espéré des fausses alertes plus la perte espérée de l'emballement, sachant qu'un emballement survient une fois tous les $N$ jours.
\end{interviewq}
